% =====================================================================
%  مقدمة الناظم (١): الأبيات ١–١٩
%  البسملة والصلاة والحمد، ثم فضل القرآن وقارئه
% =====================================================================
\chapter{مقدمة الناظم (١): البسملة والحمد وفضل القرآن}

\begin{tanbih}[عن المقدمة]
مقدمة الشاطبية (وتسمى «خطبة الكتاب») ٩٤ بيتًا. ليس فيها خلاف بين القراء، بل هي: حمد وثناء، ثم فضل القرآن وأهله، ثم التعريف بالقراء السبعة ورواتهم، ثم \textbf{الرموز والاصطلاحات} (وهي أهم جزء لفهم ما بعده)، ثم وصايا وآداب.
لذلك نكتفي في أكثرها بمعاني الكلمات والمعنى المختصر، ونتوسع في أبيات الرموز (٤٤–٦٦).
\end{tanbih}

\section{البسملة والصلاة على النبي ﷺ والحمد}

\bayt{1}{بَدَأْتُ بِبِسْمِ اللهِ فِي النَّظْمِ أَوَّلَا}{تَبَارَكَ رَحْمَانًا رَحِيمًا وَمَوْئِلَا}
\begin{mufradat}
\kalima{النظم}{الشِّعر المنظوم، والمراد هنا هذه القصيدة.}
\kalima{تبارك}{كثُر خيره وعمّت بركته.}
\kalima{رحمانًا رحيمًا}{الرحمن: المنعم بجلائل النعم وعظائمها، والرحيم: المنعم بدقائقها.}
\kalima{مَوْئِلًا}{مرجعًا وملجأً يُلجأ إليه.}
\end{mufradat}
\maana بدأتُ قصيدتي بالبسملة أولًا، تبركًا باسم الله الذي كثرت بركته، فهو الرحمن الرحيم، وهو ملجأ اللاجئين.

\bayt{2}{وَثَنَّيْتُ صَلَّى اللهُ رَبِّي عَلَى الرِّضَا}{مُحَمَّدٍ الْمُهْدَى إِلَى النَّاسِ مُرْسَلَا}
\begin{mufradat}
\kalima{ثنّيت}{جعلتُه الثاني بعد البسملة.}
\kalima{الرِّضا}{المَرْضِيّ، أي الذي ارتضاه الله للنبوة.}
\kalima{المُهْدَى}{المُرسَل هديةً ورحمةً للناس.}
\end{mufradat}
\maana ثم ثنّيت بالصلاة على النبي محمد ﷺ الذي ارتضاه الله، وبعثه هدية لعباده، وواسطةً بينهم وبين خالقهم.

\bayt{3}{وَعِتْرَتِهِ ثُمَّ الصَّحَابَةِ ثُمَّ مَنْ}{تَلَاهُمْ عَلَى الْإِحْسَانِ بِالْخَيْرِ وُبَّلَا}
\begin{mufradat}
\kalima{عِتْرته}{أهل بيته وأقاربه الأدنَون.}
\kalima{تلاهم}{تبعهم.}
\kalima{وُبَّلًا}{جمع «وابل» وهو المطر الغزير.}
\end{mufradat}
\maana وصلى الله كذلك على آل النبي ﷺ ثم على صحابته ثم على من تبعهم بإحسان، صلاةً كثيرة تنزل عليهم بالخير كالمطر الغزير.

\bayt{4}{وَثَلَّثْتُ أَنَّ الْحَمْدَ لِلَّهِ دَائِمًا}{وَمَا لَيْسَ مَبْدُوءًا بِهِ أَجْذَمُ الْعَلَا}
\begin{mufradat}
\kalima{ثلّثت}{جعلتُه الثالث.}
\kalima{أجذم العلا}{مقطوع الرِّفعة والبركة. والأجذم: المقطوع.}
\end{mufradat}
\maana وجعلت الثالث: الحمد لله دائمًا؛ لأن كل أمر مهم لا يُبدأ بحمد الله فهو ناقص البركة. وفيه إشارة إلى الحديث: «كل أمر ذي بال لا يُبدأ فيه بالحمد لله فهو أجذم».

\section{فضل القرآن}

\bayt{5}{وَبَعْدُ فَحَبْلُ اللهِ فِينَا كِتَابُهُ}{فَجَاهِدْ بِهِ حِبْلَ الْعِدَا مُتَحَبِّلَا}
\begin{mufradat}
\kalima{حبل الله}{القرآن، سُمّي حبلًا لأنه سبب النجاة، والمتمسك به لا يضل.}
\kalima{حِبْل العِدا}{كيد الأعداء ومكرهم (الحِبْل بالكسر: الداهية).}
\kalima{متحبِّلًا}{متمسكًا به كمن يمسك بالحبل.}
\end{mufradat}
\maana أما بعد: فالقرآن هو حبل الله الممدود بيننا، فجاهد به كيد أعدائك متمسكًا به. وفيه إشارة إلى قوله تعالى: \ayah{وَٱعۡتَصِمُواْ بِحَبۡلِ ٱللَّهِ جَمِيعࣰا} \src{آل عمران: ١٠٣}.

\bayt{6}{وَأَخْلِقْ بِهِ إِذْ لَيْسَ يَخْلُقُ جِدَّةً}{جَدِيدًا مُوَالِيهِ عَلَى الْجِدِّ مُقْبِلَا}
\begin{mufradat}
\kalima{أَخْلِقْ به}{ما أجدره وما أحقّه (صيغة تعجب).}
\kalima{يَخْلُق}{يَبْلَى ويقدُم.}
\kalima{جِدّة}{الجِدَّة: كون الشيء جديدًا.}
\kalima{مواليه}{ملازمه ومحبه.}
\end{mufradat}
\maana ما أحقّ القرآن بأن يُتمسك به! فهو لا يبلى ولا تذهب جِدّته مهما كثر ترداده، وصاحبه الملازم له يبقى متجددًا مُقبِلًا على الجِدّ والاجتهاد.

\bayt{7}{وَقَارِئُهُ الْمَرْضِيُّ قَرَّ مِثَالُهُ}{كَالْأُتْرُجِّ حَالَيْهِ مُرِيحًا وَمُوكَلَا}
\begin{mufradat}
\kalima{قرّ مثاله}{ثبت مَثَلُه وصفته.}
\kalima{الأُتْرُجّ}{ثمرة طيبة الطعم والرائحة.}
\kalima{حاليه}{في حالتيه.}
\kalima{مُريحًا ومُوكَلًا}{ذا ريح طيبة، ومأكولًا طيب الطعم.}
\end{mufradat}
\maana قارئ القرآن العامل به الذي يرضاه الله ثبت مَثَلُه في الحديث بأنه كالأترجة: طيبة في حالتيها، ريحها طيب وطعمها طيب. وهذا من قوله ﷺ: «مثل المؤمن الذي يقرأ القرآن كمثل الأترجة: ريحها طيب وطعمها طيب» (متفق عليه).

\bayt{8}{هُوَ الْمُرْتَضَى أَمًّا إِذَا كَانَ أُمَّةً}{وَيَمَّمَهُ ظِلُّ الرَّزَانَةِ قَنْقَلَا}
\begin{mufradat}
\kalima{أَمًّا}{مقصودًا يُقتدى به (من أَمَّ يؤُمّ: قصد).}
\kalima{أُمّة}{رجلًا جامعًا لخصال الخير، كقوله تعالى: \ayah{إِنَّ إِبۡرَٰهِيمَ كَانَ أُمَّةࣰ} \src{النحل: ١٢٠}.}
\kalima{يمّمه}{قصده.}
\kalima{الرزانة}{الوقار والثبات.}
\kalima{قَنْقَلًا}{كثيرًا وافرًا (وأصل القنقل: المكيال الضخم).}
\end{mufradat}
\maana هذا القارئ هو المرضي الذي يُقصد ويُقتدى به، إذا جمع خصال الخير وأقبل عليه الوقار والسكينة بكثرة.

\bayt{9}{هُوَ الْحُرُّ إِنْ كَانَ الْحَرِيَّ حَوَارِيًا}{لَهُ بِتَحَرِّيهِ إِلَى أَنْ تَنَبَّلَا}
\begin{mufradat}
\kalima{الحُرّ}{الكريم الذي لم يستعبده هواه.}
\kalima{الحَرِيّ}{الجدير والحقيق.}
\kalima{حواريًّا}{ناصرًا مخلصًا (كحواريّي الأنبياء).}
\kalima{تحرّيه}{طلبه الأحرى والأصوب.}
\kalima{تنبّلا}{صار نبيلًا فاضلًا.}
\end{mufradat}
\maana وهو الكريم حقًّا إن كان جديرًا بأن يكون ناصرًا مخلصًا للقرآن، يتحرّى العمل به ويقصد الصواب فيه حتى يبلغ النبل والفضل.

\bayt{10}{وَإِنَّ كِتَابَ اللهِ أَوْثَقُ شَافِعٍ}{وَأَغْنَى غَنَاءٍ وَاهِبًا مُتَفَضِّلَا}
\begin{mufradat}
\kalima{أوثق شافع}{أوثق من يشفع، أي شفاعته مقبولة.}
\kalima{أغنى غَناء}{أعظم نفعًا وكفاية.}
\kalima{واهبًا متفضلًا}{معطيًا الأجر بلا مقابل.}
\end{mufradat}
\maana القرآن أوثق شفيع لصاحبه يوم القيامة، وأعظم ما ينفعه ويغنيه، يهب له الثواب تفضلًا. وفيه إشارة إلى حديث: «اقرؤوا القرآن فإنه يأتي يوم القيامة شفيعًا لأصحابه» (رواه مسلم).

\bayt{11}{وَخَيْرُ جَلِيسٍ لَا يُمَلُّ حَدِيثُهُ}{وَتَرْدَادُهُ يَزْدَادُ فِيهِ تَجَمُّلَا}
\begin{mufradat}
\kalima{ترداده}{تكراره.}
\kalima{تجمّلًا}{حُسنًا وجمالًا.}
\end{mufradat}
\maana والقرآن خير جليس، لا يُملّ حديثه، وكلما كرّرته ازداد حسنًا وجمالًا في نفسك، بخلاف كلام البشر الذي يُملّ مع التكرار.

\bayt{12}{وَحَيْثُ الْفَتَى يَرْتَاعُ فِي ظُلُمَاتِهِ}{مِنَ الْقَبْرِ يَلْقَاهُ سَنًا مُتَهَلِّلَا}
\begin{mufradat}
\kalima{يرتاع}{يفزع ويخاف.}
\kalima{سَنًا}{ضوءًا ونورًا.}
\kalima{متهللًا}{مشرقًا مستبشرًا.}
\end{mufradat}
\maana وحين يفزع الإنسان في ظلمات القبر، يلقاه القرآن نورًا مشرقًا يؤنسه.

\bayt{13}{هُنَالِكَ يَهْنِيهِ مَقِيلًا وَرَوْضَةً}{وَمِنْ أَجْلِهِ فِي ذِرْوَةِ الْعِزِّ يَجْتَلِي}
\begin{mufradat}
\kalima{يَهْنيه}{يُهنِّئه (أصلها يهنئه بالهمز فخُففت).}
\kalima{مَقِيلًا}{مكان الراحة وقت القيلولة، والمراد القبر.}
\kalima{ذِرْوة العز}{أعلى درجات العزة.}
\kalima{يجتلي}{يظهر ويُرى في أبهى صورة.}
\end{mufradat}
\maana هناك في القبر يهنئه القرآن بمكان راحة وروضة من رياض الجنة، وبسببه يظهر صاحبه في أعلى منازل العزة.

\bayt{14}{يُنَاشِدُ فِي إِرْضَائِهِ لِحَبِيبِهِ}{وَأَجْدِرْ بِهِ سُؤْلًا إِلَيْهِ مُوَصَّلَا}
\begin{mufradat}
\kalima{يناشد}{يسأل ربه ويلحّ في السؤال.}
\kalima{لحبيبه}{أي لقارئه الذي أحبه.}
\kalima{أجدِرْ به}{ما أحقّه (تعجب).}
\kalima{سؤلًا موصّلًا}{طلبًا مقبولًا مُبلَّغًا.}
\end{mufradat}
\maana يسأل القرآن ربه أن يرضى عن قارئه، وما أحقّ هذا الطلب بالقبول. وفيه إشارة إلى حديث: «يجيء القرآن يوم القيامة فيقول: يا رب حَلِّه، فيُلبس تاج الكرامة... ثم يقول: يا رب ارضَ عنه، فيرضى عنه» (رواه الترمذي).

\bayt{15}{فَيَا أَيُّهَا الْقَارِي بِهِ مُتَمَسِّكًا}{مُجِلًّا لَهُ فِي كُلِّ حَالٍ مُبَجِّلَا}
\begin{mufradat}
\kalima{مُجِلًّا مُبجِّلًا}{معظّمًا مكرّمًا.}
\end{mufradat}
\maana يا أيها القارئ المتمسك بالقرآن العامل به، المعظم له في كل أحوالك (وجواب النداء في البيت التالي).

\bayt{16}{هَنِيئًا مَرِيئًا وَالِدَاكَ عَلَيْهِمَا}{مَلَابِسُ أَنْوَارٍ مِنَ التَّاجِ وَالْحُلَى}
\begin{mufradat}
\kalima{هنيئًا مريئًا}{دعاء بأن يكون الخير سائغًا لا تعب فيه.}
\kalima{الحُلَى}{جمع حِلية، وهي الزينة.}
\end{mufradat}
\maana هنيئًا لك ما أعدّه الله لك، حتى إن والديك يُلبَسان يوم القيامة تاجًا وحُللًا من نور. وفيه إشارة إلى حديث: «من قرأ القرآن وعمل بما فيه أُلبس والداه تاجًا يوم القيامة ضوؤه أحسن من ضوء الشمس...» (رواه أبو داود).

\bayt{17}{فَمَا ظَنُّكُمْ بِالنَّجْلِ عِنْدَ جَزَائِهِ}{أُولَئِكَ أَهْلُ اللهِ وَالصَّفْوَةُ الْمَلَا}
\begin{mufradat}
\kalima{النَّجْل}{الولد، والمراد القارئ نفسه.}
\kalima{الصَّفْوة}{الخيار.}
\kalima{المَلَا}{الأشراف.}
\end{mufradat}
\maana إذا كان هذا إكرام الوالدين، فكيف يكون إكرام الولد القارئ نفسه؟! إنهم أهل الله وخاصته، كما في الحديث: «أهل القرآن هم أهل الله وخاصته» (رواه ابن ماجه).

\bayt{18}{أُولُو الْبِرِّ وَالْإِحْسَانِ وَالصَّبْرِ وَالتُّقَى}{حُلَاهُمْ بِهَا جَاءَ الْقُرَانُ مُفَصَّلَا}
\begin{mufradat}
\kalima{حُلاهم}{صفاتهم الحسنة.}
\end{mufradat}
\maana وهم أصحاب البر والإحسان والصبر والتقوى، وهذه صفاتهم قد جاء القرآن بتفصيلها في آيات كثيرة.

\bayt{19}{عَلَيْكَ بِهَا مَا عِشْتَ فِيهَا مُنَافِسًا}{وَبِعْ نَفْسَكَ الدُّنْيَا بِأَنْفَاسِهَا الْعُلَى}
\begin{mufradat}
\kalima{عليك بها}{الزَمْ هذه الصفات.}
\kalima{منافسًا}{مسابقًا غيرك راغبًا فيها.}
\kalima{أنفاسها العُلى}{النفيس الغالي من الصفات العالية.}
\end{mufradat}
\maana الزم هذه الصفات طول حياتك متسابقًا فيها، واجعل الدنيا ثمنًا تبذله لتنال نفسُك تلك المراتب العالية، أي قدّم الآخرة على الدنيا.
